\section{Conclusion}\label{ch:conclusion}

Feasibility is a qualified yes, and the qualifications are the substance.
Shielded on-device SLM control significantly beats the fixed-time floor on
the study's most complex junction geometry and, once distilled, on a
demand-calibrated grid; the same controller is worse on a saturated
arterial whether distilled or not, and elsewhere indistinguishable from
that floor except on the oversaturated grid, where the undistilled
controller is significantly worse. The classical baseline changes sign
with the network, and what looked like uniform weakness on real geometry
was a three-seed result that replication retracted. The whole loop, from
frontier-teacher distillation through QLoRA training to int4 serving, fits
on one consumer device, and there distillation saturates: a student six
times larger lands on statistically indistinguishable decision accuracy,
pointing at the distillation ceiling rather than model capacity.

The central finding reframes what fine-tuning is \emph{for} in a shielded
controller. On the maps that could not clear, the 97\%-accurate student
does not significantly beat the floor, and neither does anything else
inside the shield's envelope, though I ran no shield ablation to isolate
that compression. Where the network \emph{can} flow, fine-tuning moves
delay itself, on every one of thirty seeds (\S\ref{sec:res-calib});
applying the same calibration to the arterial corridor did not reproduce
it, and I report that as the failed uniformity test it is. Fine-tuning
moves authorship on every topology I tested (\S\ref{sec:res-authorship}),
and the audit layer proves who decided what, so authorship is what makes
that audit concern the certified model rather than its fallback.

The nulls bound the claim, each reported with its mechanism: the corridor
harm, the stealth attacker the trust ledger misses, the privileged-teacher
ceiling, and the three serving-reliability findings. Within them, an
audited, shielded, on-device SLM junction controller is possible today, at
0.6B, on a desktop-class 6\,GB GPU, with published figures for
cabinet-class hardware suggesting but not establishing the same envelope
(Appendix~\ref{app:deploy}). For a signal engineer the reading is narrow:
deploy this where the network can still clear or where no engineered plan
exists, not on a saturated arterial; calibrate a test network to real
demand before believing any delay figure; and treat the audit trail, not
the delay, as what certification would rest on. What this study settles is
narrower and firmer: fine-tuning does not change the audit trail, which
verifies identically in every arm; it raises the share of the record the
model owns, and it moves delay only where the network can clear, and not
on every network that can.
